\textbf{Рекурсия и перебор}. Разделяй и властвуй; встреча посередине; Отсечения, A*, двунаправленный поиск; Итеративное углубление DFS / A*; Поиск характерных точек за N log N.

\textbf{Анализ сложности}. Мастер-теорема; амортизация.

\textbf{Жадные алгоритмы}. Расписания, максимальный подотрезок; Инварианты, код Хаффмана.

\textbf{Графы}. динамические графы, дополнительный учёт; BFS, DFS, деревья обхода; Дейкстра, Беллман — Форд, Флойд; Остов: Прим / Краскал; Теорема Кёнига, вершинное покрытие; Поток минимальной стоимости; Переключение Ловаса; Теорема о числе остовов; Паросочетание в общем графе; Кун, Хопкрофт — Карп, теорема Холла; Последовательности степеней; Эйлеров путь и цикл; Потоки: увеличивающие пути, Диниц; Минимальное покрытие путями; Топологическая сортировка, SCC, 2-SAT; Мосты, точки сочленения, двусвязность; Раскраска рёбер и вершин; Двудольность; деревья; Раскраска через покрытие масок за $3^n$; Диаметр, центроид, k-й кратчайший путь; Кратчайший цикл.

\textbf{Динамика}. рюкзак, размен монет; Наибольшая общая подпоследовательность; Наибольшая возрастающая подпоследовательность; Число путей / кратчайший путь в DAG; ДП по отрезкам, маскам, вероятностям, деревьям; Покрытие множества за $3^n$; Разделяй и властвуй, оптимизация Кнута; CHT, монотонная очередь; RMQ, разреженная таблица; Битонический цикл; Перебор по наиболее ограниченной части.

\textbf{Комбинаторика}. биномиальные коэффициенты, принцип Дирихле; Включения-исключения, Каталан, теорема Пика.

\textbf{Теория чисел}. целые части, делимость, Евклид; Модульные умножение, обратный, степень; КТО, Ферма, Эйлер, функция phi; Число Фробениуса, квадратичная взаимность; Поллард, Миллер — Рабин, подъём Гензеля; Прыжки Виета.

\textbf{Игры}. Комбинаторные игры, дерево игры, минимакс; Ним, числа Гранди; Игры на графах с циклами и без; Партизанские и общие игры без повторов; Альфа-бета отсечение.

\textbf{Вероятности; оптимизация}. Бинарный / тернарный поиск; Унимодальность, выпуклость; Бинарный поиск по производной.

\textbf{Численные методы}. Интегрирование, метод Ньютона; Поиск корня, золотое сечение.

\textbf{Матрицы}. Гаусс, возведение в степень.

\textbf{Сортировка}. поразрядная сортировка.

\textbf{Геометрия}. Векторы, скалярное и векторное произведения; Выпуклая оболочка, отсечение многоугольника; Ближайшая пара, сжатие координат; Квадродерево, KD-tree; Все пересечения отрезков.

\textbf{Заметающая прямая}. События и дискретизация; Поворот луча, движение прямой; Дискретные вторые производные.

\textbf{Строки}. наибольшая общая подстрока; Палиндромные подпоследовательности; КМП, бор, полиномиальные хеши; Суффиксный массив, дерево, автомат; Ахо — Корасик, Манакер; Списки позиций символов.

\textbf{Структуры данных}. LCA, двоичные подъёмы; Перенос данных в дереве, малое в большое; HLD, центроидная декомпозиция; Ленивое проталкивание, балансировка; CHT, монотонные стеки и очереди; Скользящее окно, очередь на двух стеках; Персистентное ДО.
